\sectionheading{5.25 Python Programming Language-এর পরিচিতি}

বর্তমানে শিক্ষা, Software Development, Data Analysis,
Artificial Intelligence, Web Developmentসহ বিভিন্ন ক্ষেত্রে
বহুল ব্যবহৃত Programming Language হলো \textbf{Python}।

সহজ ও বোধগম্য Syntax-এর কারণে নতুনদের Programming শেখার
ক্ষেত্রেও Python বিশেষভাবে জনপ্রিয়।

%================================================
% Definition
%================================================

\booksubheading{Python কী?}

\begin{definitionbox}

\textbf{\color{darkTeal}সংজ্ঞা}

\textbf{Python হলো একটি High Level ও General Purpose
Programming Language}, যার সাহায্যে বিভিন্ন ধরনের
Computer Program ও Software তৈরি করা যায়।

\end{definitionbox}

\textbf{সহজভাবে মনে রাখবে—}

\begin{center}

\textbf{সহজ Syntax}
\(\longrightarrow\)
\textbf{বিভিন্ন ধরনের কাজ}
\(\longrightarrow\)
\textbf{Python}

\end{center}

অর্থাৎ Python শুধু একটি নির্দিষ্ট ধরনের কাজের জন্য নয়;
বিভিন্ন ধরনের Problem সমাধান ও Software Development-এ
এটি ব্যবহার করা যায়।

%================================================
% History
%================================================

\booksubheading{Python-এর সংক্ষিপ্ত ইতিহাস}

Python Programming Language তৈরি করেন
\textbf{Guido van Rossum}।

তিনি ১৯৮০-এর দশকের শেষদিকে Python তৈরির কাজ শুরু করেন এবং
\textbf{1991 সালে Python-এর প্রথম সংস্করণ প্রকাশিত হয়}।

\medskip

\begin{tabularx}
{\linewidth}
{
|>{\bfseries}p{0.30\linewidth}
|Y|
}

\hline

\rowcolor{softBlue}

\multicolumn{1}{|c|}{\textbf{বিষয়}}
&
\multicolumn{1}{c|}{\textbf{তথ্য}}

\\

\hline

ভাষার নাম
&
Python
\\

\hline

নির্মাতা
&
\textbf{Guido van Rossum}
\\

\hline

প্রথম প্রকাশ
&
\textbf{1991}
\\

\hline

ভাষার স্তর
&
High Level Language
\\

\hline

ব্যবহারের ধরন
&
General Purpose Programming
\\

\hline

\end{tabularx}

\begin{rememberbox}

\textbf{\color{deepOrange}MCQ তথ্য}

\textbf{Python-এর নির্মাতা}
\(\rightarrow\)
Guido van Rossum

\textbf{প্রথম প্রকাশ}
\(\rightarrow\)
1991

\end{rememberbox}

%================================================
% Naming
%================================================

\booksubheading{Python নামটি কীভাবে এসেছে?}

Python নামটি \textbf{Python সাপের নাম থেকে নেওয়া হয়নি}।

Guido van Rossum ব্রিটিশ Comedy Group
\textbf{Monty Python}-এর অনুষ্ঠান পছন্দ করতেন।
সেখান থেকেই Programming Languageটির নাম
\textbf{Python} রাখা হয়।

\begin{rememberbox}

\textbf{\color{deepOrange}MCQ Trap}

\textbf{Python নামের উৎস}
\(\rightarrow\)
Monty Python

\medskip

\textbf{Python নামের উৎস}
\(\neq\)
Python সাপ

\end{rememberbox}

%================================================
% Type of Language
%================================================

\booksubheading{Python কী ধরনের Programming Language?}

Python-কে সাধারণভাবে—

\begin{itemize}

\item \textbf{High Level Language}

\item \textbf{General Purpose Programming Language}

\item \textbf{Interpreted Language}

\end{itemize}

হিসেবে ব্যবহার করা হয়।

\begin{tabularx}
{\linewidth}
{
|>{\bfseries}p{0.30\linewidth}
|Y|
}

\hline

\rowcolor{softBlue}

\multicolumn{1}{|c|}{\textbf{ধরন}}
&
\multicolumn{1}{c|}{\textbf{সহজ ধারণা}}

\\

\hline

High Level
&
মানুষের জন্য তুলনামূলক সহজে বোঝা ও লেখা যায়।
\\

\hline

General Purpose
&
বিভিন্ন ধরনের Program ও Software তৈরিতে ব্যবহার করা যায়।
\\

\hline

Interpreted
&
Python Program Interpreter-এর মাধ্যমে Execute করা হয়।
\\

\hline

\end{tabularx}

\begin{rememberbox}

এই বৈশিষ্ট্যগুলো পরবর্তী টপিক
\textbf{5.26 Python Programming Language-এর বৈশিষ্ট্য}
অংশে বিস্তারিত আলোচনা করা হবে।

\end{rememberbox}

%================================================
% Uses
%================================================

\booksubheading{Python কোথায় ব্যবহার করা হয়?}

Python বিভিন্ন ধরনের কাজে ব্যবহার করা হয়। যেমন—

\begin{tabularx}
{\linewidth}
{
|>{\bfseries}p{0.36\linewidth}
|Y|
}

\hline

\rowcolor{softBlue}

\multicolumn{1}{|c|}{\textbf{ক্ষেত্র}}
&
\multicolumn{1}{c|}{\textbf{ব্যবহার}}

\\

\hline

Web Development
&
Website ও Web Application তৈরি
\\

\hline

Data Analysis
&
Data বিশ্লেষণ
\\

\hline

Artificial Intelligence
&
AI ভিত্তিক Application তৈরি
\\

\hline

Machine Learning
&
Machine Learning Model তৈরি ও ব্যবহার
\\

\hline

Automation
&
পুনরাবৃত্তিমূলক কাজ স্বয়ংক্রিয় করা
\\

\hline

Scientific Computing
&
বৈজ্ঞানিক ও সংখ্যাভিত্তিক হিসাব
\\

\hline

Application Development
&
বিভিন্ন ধরনের Software/Application তৈরি
\\

\hline

Education
&
Programming শেখানো ও Problem Solving
\\

\hline

\end{tabularx}

\begin{rememberbox}

এখানে প্রতিটি ক্ষেত্রের বিস্তারিত শেখার প্রয়োজন নেই।
\textbf{Python বিভিন্ন ধরনের কাজে ব্যবহৃত হয়}—
এই মূল ধারণাটি মনে রাখবে।

\end{rememberbox}

%================================================
% Why Python in HSC
%================================================

\booksubheading{Python শেখার মাধ্যমে কী করা যাবে?}

Python শেখার মাধ্যমে শিক্ষার্থী—

\begin{compacttwocolpoints}

\item Programming-এর মৌলিক ধারণা বুঝতে পারবে।

\item Variable ও Data Type ব্যবহার করতে পারবে।

\item Input ও Output নিয়ে কাজ করতে পারবে।

\item Operator ব্যবহার করে Calculation করতে পারবে।

\item Condition ব্যবহার করে সিদ্ধান্তভিত্তিক Program লিখতে পারবে।

\item Loop ব্যবহার করে পুনরাবৃত্তিমূলক কাজ করতে পারবে।

\item Function ব্যবহার করতে পারবে।

\item মৌলিক Data Structure নিয়ে কাজ করতে পারবে।

\item Algorithm ও Flowchart-কে Program Code-এ রূপ দিতে পারবে।

\end{compacttwocolpoints}

%================================================
% Algorithm to Python
%================================================

\booksubheading{Algorithm থেকে Python Program}

আগের অংশে আমরা Problem সমাধানের Algorithm ও Flowchart
সম্পর্কে জেনেছি।

এখন সেই Solution-কে Python Program-এ রূপ দেওয়া হবে।

\begin{center}
\textbf{Problem}
\(\longrightarrow\)
\textbf{Algorithm}
\(\longrightarrow\)
\textbf{Flowchart}
\(\longrightarrow\)
\textbf{Python Code}
\(\longrightarrow\)
\textbf{Output}
\end{center}

\begin{rememberbox}

\textbf{\color{deepOrange}মনে রাখবে}

Algorithm ও Flowchart হলো
\textbf{Solution-এর পরিকল্পনা ও চিত্রভিত্তিক প্রকাশ}।

Python Code হলো সেই Solution-কে
\textbf{Programming Language-এ বাস্তবায়ন}।

\end{rememberbox}

%================================================
% First Python Program
%================================================

\booksubheading{একটি ছোট Python Program}

Python Code দেখতে কেমন তা বোঝার জন্য একটি ছোট উদাহরণ—

\begin{examplebox}{উদাহরণ: Hello, World! প্রদর্শন}

\begin{codebox}

\begin{english}
\ttfamily
print("Hello, World!")
\end{english}

\end{codebox}

এখানে
\textenglish{\ttfamily print()}
ব্যবহার করে Screen-এ

\textenglish{\ttfamily Hello, World!}

Output প্রদর্শন করা হয়েছে।

\end{examplebox}

\begin{rememberbox}

\textbf{\color{deepOrange}এখন শুধু ধারণা নাও}

এখনই
\textenglish{\ttfamily print()}
-এর সম্পূর্ণ Syntax মুখস্থ করার প্রয়োজন নেই।

\textbf{5.35 Input ও Output}
অংশে এটি বিস্তারিতভাবে আলোচনা করা হবে।

\end{rememberbox}

%================================================
% Python 3
%================================================

\booksubheading{Python 2 ও Python 3}

Python-এর বিভিন্ন Version প্রকাশিত হয়েছে।

বর্তমানে Programming শেখা ও নতুন Program তৈরিতে
\textbf{Python 3} ব্যাপকভাবে ব্যবহৃত হয়।

এই বইয়ের পরবর্তী Python Code ও উদাহরণগুলো
\textbf{Python 3 Syntax} অনুসরণ করে লেখা হবে।

\begin{rememberbox}

\textbf{\color{deepOrange}মনে রাখবে}

এই বইয়ের Python Program
\(\rightarrow\)
\textbf{Python 3}

\end{rememberbox}

%================================================
% Common Mistakes
%================================================

\booksubheading{সাধারণ ভুল}

\begin{compacttwocolpoints}

\item Python-এর নির্মাতা হিসেবে অন্য কারও নাম লেখা।

\item Python নামটি সাপ থেকে এসেছে মনে করা।

\item Python-কে Low Level Language বলা।

\item Python-কে শুধু AI বা শুধু Web Development-এর
Language মনে করা।

\item Algorithm, Flowchart ও Python Code-কে একই বিষয় মনে করা।

\item Python 2 ও Python 3-এর Syntax একসঙ্গে মিশিয়ে ব্যবহার করা।

\end{compacttwocolpoints}

%================================================
% Final Reminder
%================================================

\begin{rememberbox}

\textbf{\color{deepOrange}মনে রাখবে}

\begin{itemize}

\item
\textbf{Python} একটি High Level ও
General Purpose Programming Language।

\item
Python তৈরি করেন
\textbf{Guido van Rossum}।

\item
Python-এর প্রথম সংস্করণ প্রকাশিত হয়
\textbf{1991 সালে}।

\item
Python নামটি
\textbf{Monty Python}
থেকে নেওয়া হয়েছে।

\item
Python বিভিন্ন ধরনের Software ও Problem Solving-এ
ব্যবহার করা যায়।

\item
Web Development, Data Analysis, AI,
Machine Learning ও Automationসহ বিভিন্ন ক্ষেত্রে
Python ব্যবহৃত হয়।

\item
এই বইয়ের Programগুলো
\textbf{Python 3}
অনুসরণ করবে।

\item
Algorithm ও Flowchart-এর Solution-কে
Python Code-এ বাস্তবায়ন করা যায়।

\item
দ্রুত মনে রাখবে—

\begin{center}

\textbf{Guido van Rossum}
\(\longrightarrow\)
\textbf{Python}
\(\longrightarrow\)
\textbf{1991}

\end{center}

\item
Programming-এর ধারাবাহিকতা—

\begin{center}

\textbf{Problem}
\(\longrightarrow\)
\textbf{Algorithm}
\(\longrightarrow\)
\textbf{Flowchart}
\(\longrightarrow\)
\textbf{Python Code}
\(\longrightarrow\)
\textbf{Output}

\end{center}

\end{itemize}

\end{rememberbox}

\small

%================================================
% 5.26 Python Features
%================================================

